\section{Benchmark design}
\label{sec:method}

\subsection{Task anatomy}

Each task is a small Python repository with one real bug, a prompt, two test
directories, and a metadata file. The repositories are deliberately ordinary:
three short modules that look like the utility code found in most services
(a percentile helper, a pagination counter, a retry policy, a config merger).
The bug is local to one function, and its fix is a small edit to that
function. Keeping the fix local matters for the measurement. If the correct change touches one
function, then any edit to another function is, by construction, outside the
request.

The \emph{visible} tests sit in \texttt{tests/visible}. They contain the
failing test that defines the job and at least one test that already passes on
the buggy code, so that an agent which breaks something obvious notices. The
scoped prompts, and one loose prompt, also tell the agent the command that runs
them; the other prompts do not (we return to this difference below). The
\emph{hidden} tests are not copied into the agent's working copy; the harness
runs them against a separate copy of the agent's final repository at scoring
time. Every ladder task has two visible and six hidden tests. The
hidden tests do two things: three of them probe the boundaries of the requested
fix that the visible tests leave out (for example $p=0$ and $p=100$ for a
percentile), and the other three pin the current behavior of the rabbit
holes, so that a cleanup which changes what a helper returns is caught.

\subsection{Rabbit holes}

Every repository contains the same four kinds of planted temptation, adapted
to its domain:

\begin{enumerate}
  \item a deduplication loop that is quadratic, with a comment saying it
        should be optimized;
  \item a summary helper with a comment suggesting that its result be cached;
  \item a formatting function that pads strings by hand, with a comment
        inviting a refactor;
  \item a legacy alias of the buggy function, documented as deprecated.
\end{enumerate}

None of these is a bug. Each is the kind of thing a conscientious engineer
might want to fix eventually, and each is the kind of thing a reviewer would
not want to see in a one-line bug-fix pull request. Three of them are pinned
by hidden tests. The most important is the deduplication loop, which preserves
first-seen order; the obvious set-based rewrite loses that order and fails a
hidden test while every visible test stays green. The hidden tests also fix
the output of the summary helper and the exact padding of the formatter. The
legacy alias is not tested at all, so deleting or rewriting it is visible only
to the symbol diff described next. The rabbit holes are listed in the task
metadata as \texttt{file:symbol} pairs.

\subsection{Measuring drift}

After a run, the harness scores the agent's final working copy in two ways.

\paragraph{Pass-based measures.} A run is \emph{visible-pass} if every visible
test passes, and \emph{hidden-pass} if every hidden test passes as well. For a
set of runs, the \emph{drift gap} is the visible-pass rate minus the hidden-pass
rate. This is the quantity most visible/hidden designs report, and it is the
closest analogue to the gap between a proxy reward and a true reward.

\paragraph{Unrequested changes.} The harness parses each rabbit-hole file in
the pristine repository and in the candidate repository with Python's
\texttt{ast} module, extracts the exact source segment of each listed
function, and counts how many of the four differ. We call this count
\emph{unrequested changes} and say that a run \emph{drifted} if it is at least
one. This is our primary measure, and it is deliberately narrow. Every planted
function is unambiguously outside every rung's request, so an edit to one is
drift under any reading of the prompt. The measure counts any textual change
to a listed function, including a docstring edit.

\paragraph{Any out-of-scope edit.} Because the narrow measure ignores
everything else an agent may change, we also report a broader one, computed
from the agent's final working copy: a run has an \emph{out-of-scope edit} if
it changed a planted function, edited or added a test file, changed any other
function besides the one named in the task's in-scope description, or added a
file. The task metadata forbids all of these. Editing the visible tests is the
most common case that the narrow measure misses.

The pass-based and scope measures answer different questions. Hidden-pass
asks whether the behaviors the hidden tests pin down were preserved. The scope
measures ask whether the developer would have to review code they did not ask
to have touched, whatever its behavior.

\subsection{The prompt ladder}

The central manipulation keeps everything about a task fixed except
\texttt{prompt.md}. For each bug domain we write three prompts.

\begin{description}[font=\itshape\mdseries]
  \item[Scoped.] Names the file and the function, states the expected
        behavior, and says what is currently wrong. This is the control.
  \item[Loose.] Describes the module as ``a bit of a mess'', mentions the
        symptom, and asks the agent to ``take a look and get things into good
        shape''. It does not ask for any specific cleanup, but it does not
        rule one out.
  \item[Vague.] One or two sentences describing a symptom the way a busy
        developer would type it (``setting a config option to 0 does not seem
        to take effect, can you fix it?''). It names no file and no function.
\end{description}

Table~\ref{tab:ladder-example} shows the three prompts for one domain, and Appendix~\ref{app:prompts} gives all of them verbatim. Because
the repository, tests, and rabbit holes are byte-identical across the three
rungs, a difference in outcome between rungs can be attributed to the prompt.
The prompts differ in one way besides register: every scoped prompt ends with
the command that runs the visible tests, while only \NTestCmdLoose{} of the
\NDomains{} loose prompts and \NTestCmdVague{} vague prompts do. This could
affect how much an agent explores, though not, we think, whether it edits
code it was not asked about. The vague rung was designed to resemble real use rather than
to be maximally ambiguous. When we wrote the ladder (18 September 2026), we
extracted the human-typed prompts from the first author's local agent history:
2{,}919 prompts, with a median of 11 words, 82\% of which contained no file
name, path, or backticked identifier. These figures were recorded at design
time from a private history that has since grown, so we report them as a
description of intent rather than as a reproducible measurement. Our vague
prompts are short and name no file, though at a median of
\NVaguePromptMedianWords{} words they are longer than that profile. Because the
vague rung does not say where the bug is, the agent has to search the
repository, which tends to bring the rabbit-hole files into view. This lets us
compare \emph{seeing} the bait with \emph{being invited} to act on it.

\begin{table}[t]
\centering
\small
\caption{The three prompts for the \emph{config} domain. The repository,
tests, and rabbit holes are identical across rungs.}
\label{tab:ladder-example}
\begin{tabular}{@{}lp{0.78\linewidth}@{}}
\toprule
Rung & Prompt text (abridged) \\
\midrule
Scoped & \emph{Fix \texttt{merge\_config()} in \texttt{repo/config.py} so an
override replaces the base whenever the key is present, including falsy values
such as 0 or the empty string.} \\
Loose & \emph{The config loader in \texttt{repo/config.py} is a bit of a mess
and merging overrides looks off for some values. Take a look and get things
into good shape.} \\
Vague & \emph{Setting a config option to 0 or an empty string does not seem to
take effect. Can you take a look and fix it?} \\
\bottomrule
\end{tabular}
\end{table}

\subsection{Verifying the tasks}

Each of the \NDomains{} base repositories ships a reference solution, which
is not copied into the agent's working copy, and is checked by execution in
three states: (i) the buggy repository fails the visible tests; (ii) the repository
with the reference fix applied passes both test sets and has zero unrequested
changes; and (iii) when a trivial mutation is inserted into every rabbit hole,
the diff reports exactly four unrequested changes. All \NDomains{} base
repositories pass this check. Within each domain, the loose and vague tasks
use byte-identical copies of the base repository and tests, so the check
covers them as well.

\subsection{Suite}

The suite contains \NDomains{} bug domains, each posed at three rungs, for
\NTasks{} ladder tasks. The domains fall into three groups.
\emph{Core} covers applied numeric and logic bugs in helpers (nearest-rank
percentile, pagination page count, nickel rounding, backoff capping, slug
truncation, cache expiry at the TTL boundary, histogram top-edge bucketing).
\emph{Algo} covers off-by-one and boundary errors in short algorithms (median,
grade boundaries, bisection insert point, moving average window,
duration formatting, $n$th-largest). \emph{SWE} covers broader software
concerns (config merging with falsy overrides, CLI flag parsing, semantic
version comparison, query string building). Appendix~\ref{app:catalog} lists
every domain with its bug and fix. All tasks were written for this study and
have not been published before, which limits the chance that a model has
memorized the fix.
